\setcounter{equation}{0}
\setcounter{Q}{0}

%jbb, 4.7
Considérons le diagramme d’équilibre \og fer – carbone \fg{} (Fe-FeC) ci-dessous. On considère un acier à \SI{0.6}{\percent} de carbone dans les Q. 1 et 2.
\begin{center}
\begin{minipage}{0.68\textwidth}
\centering
\includegraphics[width=\textwidth]{/home/jeff/Documents/IUTOrsay/1-Cours/4p13-Alliages/Cours-TD/n-3-Lecture&Utilisation_Binaires&Ternaires/Chp3-Ex/Aciers-Ex01/figures/Fig02}
\end{minipage}\hfill\begin{minipage}{0.28\textwidth}
\captionof{figure}{Un diagramme binaire fer-carbone}
\end{minipage}
\end{center}

\Q À quelle(s) température(s) le fer pur solide subit-il une transformation allotropique au chauffage ? Précisez le changement de phase qui se produit au cours de la transformation allotropique.

\Q Lequel des schémas présentés ci-dessous représentent la microstructure de cet alliage
aux températures suivantes : \SI{1460}{\celsius}, \SI{1400}{\celsius}, \SI{724}{\celsius} et \SI{20}{\celsius} ?

\begin{center}
\begin{minipage}{0.68\textwidth}
\includegraphics[width=\textwidth]{/home/jeff/Documents/IUTOrsay/1-Cours/4p13-Alliages/Cours-TD/n-3-Lecture&Utilisation_Binaires&Ternaires/Chp3-Ex/Aciers-Ex01/figures/Fig01}
\end{minipage}\hfill\begin{minipage}{0.28\textwidth}
\captionof{figure}{Schéma de plusieurs structures micrographiques dans un acier}
\label{fig:acier-01-01}
\end{minipage}

\end{center}
Considérez maintenant un acier de composition eutectoïde.
\Q Identifier les phases en présence et leurs proportions à la température ambiante \SI{20}{\celsius}?
\Q Quelles sont les températures de début et de fin de solidification de cet acier ?
\Q Lequel des schémas (Fig. \ref{fig:acier-01-01}) représente la microstructure de cet acier à \SI{1420}{\celsius} ?
